\uuid{Hv2R}
\titre{Le problème du gradient vanishing}
\niveau{L3}
\module{Apprentissage automatique}
\chapitre{Réseaux de neurones récurrents}
\sousChapitre{Gradient vanishing et exploding}
\theme{Règle de dérivation en chaîne, saturation de tanh, dépendances à long terme}
\auteur{Maxime Nguyen}
\datecreate{2026-05-13}
\organisation{AMSCC}
\difficulte{3}

\contenu{
	\texte{
		On considère un RNN simple scalaire sans entrée, pour isoler la dynamique de l'état caché :
		\[
		h_t = \tanh(w_h \, h_{t-1})
		\]
		On s'intéresse à la propagation du gradient sur $T$ pas de temps. Par la règle de dérivation en chaîne :
		\[
		\frac{\partial h_T}{\partial h_0} = \prod_{t=1}^{T} \frac{\partial h_t}{\partial h_{t-1}}
		= \prod_{t=1}^{T} w_h \cdot (1 - h_t^2)
		\]
	}
	\begin{enumerate}
		\item
		\question{Justifier la formule ci-dessus en appliquant la règle de dérivation en chaîne à $h_T = f(h_{T-1}) = f(f(\cdots f(h_0)\cdots))$.}
		\reponse{
			La fonction $f : x \mapsto \tanh(w_h x)$ est composée $T$ fois. Par la règle de dérivation des fonctions composées :
			\[
			\frac{\partial h_T}{\partial h_0}
			= \frac{\partial h_T}{\partial h_{T-1}} \cdot \frac{\partial h_{T-1}}{\partial h_{T-2}} \cdots \frac{\partial h_1}{\partial h_0}
			= \prod_{t=1}^{T} f'(h_{t-1})
			\]
			Or $f'(x) = w_h \cdot (1 - \tanh(w_h x)^2) = w_h \cdot (1 - h_t^2)$, d'où la formule.
		}
		
		\item
		\question{Montrer que pour tout $h \in \mathbb{R}$, on a $0 < 1 - \tanh(h)^2 \leq 1$, avec égalité uniquement en $h = 0$. Que peut-on en conclure sur chaque facteur $w_h \cdot (1 - h_t^2)$ lorsque $|w_h| \leq 1$ ?}
		\reponse{
			On a $\tanh(h)^2 \in [0, 1)$ pour tout $h \in \mathbb{R}$ (puisque $|\tanh(h)| < 1$), donc $1 - \tanh(h)^2 \in (0, 1]$, avec la valeur $1$ atteinte seulement en $h = 0$ (car $\tanh(0) = 0$).
			
			Si $|w_h| \leq 1$, alors $|w_h \cdot (1 - h_t^2)| \leq 1$, et dès que $h_t \neq 0$, la majoration est stricte. Chaque facteur du produit est donc strictement inférieur à $1$ en valeur absolue, ce qui implique que le produit de $T$ tels facteurs tend vers $0$ exponentiellement vite quand $T \to +\infty$.
		}
		
		\item
		\question{On suppose $h_t \approx 0.9$ pour tout $t$, et $w_h = 0.8$. Calculer numériquement $\left|\dfrac{\partial h_T}{\partial h_0}\right|$ pour $T = 5$, $T = 10$, $T = 20$. Commenter.}
		\reponse{
			Chaque facteur vaut $|w_h \cdot (1 - h_t^2)| = 0.8 \times (1 - 0.9^2) = 0.8 \times (1 - 0.81) = 0.8 \times 0.19 = 0.152$.
			
			\[
			T = 5  : \quad 0.152^5 \approx 8.15 \times 10^{-5}
			\]
			\[
			T = 10 : \quad 0.152^{10} \approx 6.65 \times 10^{-9}
			\]
			\[
			T = 20 : \quad 0.152^{20} \approx 4.42 \times 10^{-17}
			\]
			Le gradient devient infinitésimal (pratiquement nul) dès les premiers pas de temps : toute information sur $h_0$ est perdue lors de la rétropropagation, rendant l'apprentissage de dépendances longues impossible.
		}
		
		\item
		\question{Montrer que le gradient ne s'évanouit pas si et seulement si $|w_h| \cdot (1 - h_t^2) \geq 1$ pour tout $t$. Quelle condition sur $w_h$ cela impose-t-il lorsque $h_t \approx 0$ ? Quel nouveau problème apparaît alors ?}
		\reponse{
			Pour que $\left|\dfrac{\partial h_T}{\partial h_0}\right|$ ne tende pas vers $0$, il faut que le produit ne s'effondre pas, ce qui requiert que chaque facteur soit au moins $1$ en valeur absolue : $|w_h| \cdot (1 - h_t^2) \geq 1$.
			
			Lorsque $h_t \approx 0$, on a $1 - h_t^2 \approx 1$, donc la condition devient $|w_h| \geq 1$.
			
			Mais si $|w_h| > 1$ et que les activations ne saturent pas, les facteurs sont tous strictement supérieurs à $1$, et le produit diverge exponentiellement : c'est le problème du \textit{gradient exploding}, qui rend l'entraînement instable (mises à jour de poids de taille arbitrairement grande).
		}
		
		\item
		\question{On cherche à minimiser une perte $\mathcal{L}$ dépendant de $h_T$. Exprimer $\dfrac{\partial \mathcal{L}}{\partial w_h}\big|_{t=1}$, la contribution du tout premier pas de temps au gradient de $w_h$, en fonction de $\dfrac{\partial \mathcal{L}}{\partial h_T}$ et du produit ci-dessus. Pourquoi les dépendances à long terme sont-elles si difficiles à apprendre ?}
		\reponse{
			Par la règle de dérivation en chaîne, la contribution du pas $t = 1$ est :
			\[
			\left.\frac{\partial \mathcal{L}}{\partial w_h}\right|_{t=1}
			= \frac{\partial \mathcal{L}}{\partial h_T} \cdot \frac{\partial h_T}{\partial h_1} \cdot \frac{\partial h_1}{\partial w_h}
			= \frac{\partial \mathcal{L}}{\partial h_T} \cdot \left(\prod_{t=2}^{T} w_h(1-h_t^2)\right) \cdot h_0(1-h_1^2)
			\]
			Ce terme est proportionnel au produit de $(T-1)$ facteurs, chacun inférieur à $1$ si $|w_h| \leq 1$. Il devient donc exponentiellement petit : le signal de gradient provenant de la perte à $T$ ne parvient pratiquement plus au premier pas de temps, empêchant le réseau d'ajuster $w_h$ pour capturer une dépendance entre $h_1$ et $h_T$.
		}
	\end{enumerate}
}